\documentclass[12pt]{article}
\usepackage{amsmath}
\usepackage{graphicx}


\addtolength{\textheight}{0.5in}
\addtolength{\topmargin}{-0.2in}
\addtolength{\evensidemargin}{-0.0in}
\addtolength{\textwidth}{0.9in}
%\raggedbottom


%\def\sst{\scriptscriptstyle}
\renewcommand{\baselinestretch}{1.1}
\renewcommand{\arraystretch}{1.2}
\parindent=20pt

\renewcommand{\topfraction}{0.9}
\renewcommand{\textfraction}{0.1}
\renewcommand{\floatpagefraction}{0.9}


\setlength{\unitlength}{1cm}

\begin{document}

\title{Test of Gaussian folding routing}
\author{Discussion file: FAU -- MSU}
\date{Started: 11. June 2024}
\maketitle


The file \texttt{Gauss\_folding.f90} provides a routine for the
Gaussian folding
\begin{equation}
  \varphi'(\mathbf{r})
  =
  \Delta r^3\pi^{3/2}\int d\mathbf{r}_2
  \exp\left(-\frac{(\mathbf{r}-\mathbf{r}_2)^2}{\Delta r^2}\right)
  \varphi(\mathbf{r}_2)
\end{equation}
acting on a s.p. wavefunction $\varphi$. It is specified for 1D
spherical calculations on a coordinate-space grid, thus serving as a
plugin to the Erlangen 1D SHF code.

The new piece of code needs yet testing. This can be done by using
oscillator functions as probes. For simplicity, we those whicn can be
expressed by monomials
\begin{equation}
  \varphi
  =
  r^lY_{ll}(\Omega)\exp\left(-\frac{r^2}{\delta r^2}\right)
  \quad.
\label{eq:initHO}
\end{equation}
The task is now to find an analytical expression for the corresponding
$\varphi'$. 

First, we express $\varphi$ in terms of derivatives. To that end, we
use the gradient formula 
\begin{eqnarray}
  \nabla_+r^lY_{ll}e^{-r^2/\delta r^2}
  &=&
  Y_{l+1,l+1}\sqrt{\frac{l+1}{2l+3}}\left(\partial_r-\frac{l}{r}\right)
  r^le^{-r^2/\delta r^2}
\nonumber\\
  &=&
  -\frac{2}{\delta r^2}\sqrt{\frac{l+1}{2l+3}}Y_{l+1,l+1}
  r^{l+1}e^{-r^2/\delta r^2}
\\
  \nabla_+
  &=&
  \partial_x+\mathrm{i}\partial_y  
\end{eqnarray}
from which we can deduce
\begin{equation}
  r^lY_{ll}(\Omega)\exp\left(-\frac{r^2}{\delta r^2}\right)
  =
  \left(\frac{-\delta r^2}{2}\right)^l
  \sqrt{\frac{(2l+1)!!}{l!}}\,Y_{00}\,
  \nabla_+^l\exp\left(-\frac{r^2}{\delta r^2}\right)
  \quad.
\label{eq:rpowerlfunction}
\end{equation}

Second, we use partial integration to shovel the derivative from
$\varphi$ to $\varphi'$. This reads
\begin{eqnarray*}
  \varphi'
  &=&
  \Delta r^3\pi^{3/2}\int d\mathbf{r}_2
  \exp\left(-\frac{(\mathbf{r}-\mathbf{r}_2)^2}{\Delta r^2}\right)
  r_2^lY_{ll}(\Omega_2)\exp\left(-\frac{r_2^2}{\delta r^2}\right)
\\
  &=&
  \Delta r^3\pi^{3/2}
  \left(\frac{-\delta r^2}{2}\right)^l
  \sqrt{\frac{(2l+1)!!}{l!}}\,Y_{00}\,
  \int d\mathbf{r}_2
  \exp\left(-\frac{(\mathbf{r}-\mathbf{r}_2)^2}{\Delta r^2}\right)
  \nabla_{2,+}^l\exp\left(-\frac{r^2}{\delta r^2}\right)
\\
  &=&
  \Delta r^3\pi^{3/2}
  \left(\frac{-\delta r^2}{2}\right)^l
  \sqrt{\frac{(2l+1)!!}{l!}}\,Y_{00}\,
  (-)^l\nabla_+^l
  \int d\mathbf{r}_2
  \exp\left(-\frac{(\mathbf{r}-\mathbf{r}_2)^2}{\Delta r^2}\right)
  \exp\left(-\frac{r^2}{\delta r^2}\right)
\\ 
  &=&
  \frac{(\Delta r^2+\delta r^2)^{3/2}}{\delta r^3}
  \left(\frac{-\delta r^2}{2}\right)^l
  \sqrt{\frac{(2l+1)!!}{l!}}\,Y_{00}\,
  (-)^l\nabla_+^l
  \exp\left(-\frac{r^2}{\Delta r^2+\delta r^2}\right)
\end{eqnarray*}
from which we conclude
\begin{equation}
  \varphi'
  =
  (-)^l\frac{(\Delta r^2+\delta r^2)^{3/2}}{\delta r^3}
  \left(\frac{\delta r^2}{\Delta r^2+\delta r^2}\right)^l
  r^lY_{ll}\exp\left(-\frac{r^2}{\Delta r^2+\delta r^2}\right)
\end{equation}
for the wavefunction $\varphi$ as given in eq. (\ref{eq:initHO}).


\end{document}
